% from ega0-22-fr.tex
\oldpage[0\textsubscript{IV-1}]{192}

Pour montrer que cet anneau local est une $L$-algèbre formellement lisse,
il suffit donc de prouver que $\widehat A_\mathfrak p\otimes_A K$ est une
$K$-algèbre formellement lisse (19.3.5, (iii) et (iv))~; d'ailleurs
$\widehat A_\mathfrak p\otimes_A K=\widehat A_\mathfrak p$ puisque
$\mathfrak p\cap A=0$. On va appliquer les critères (22.2.2) et (22.2.8).
En premier lieu, comme $A$ est régulier, il en est de même de $\widehat A$
(17.1.5), donc de $\widehat A_\mathfrak p$ (17.3.2), et la condition (i) de
(22.2.2) est vérifiée (17.1.1). Il suffit par suite (lorsque $p>0$) de
vérifier la condition b) de (22.2.8). Or, (22.3.1) montre que, pour tout
$\alpha$, $\widehat A$ est une $A$-algèbre formellement lisse relativement
à $A_\alpha$ (pour les topologies discrètes)~; par suite
\[
  \widehat A\otimes_{A_\alpha}K_\alpha
  =\widehat A\otimes_A(A\otimes_{A_\alpha}K_\alpha)
\]
est une $(A\otimes_{A_\alpha}K_\alpha)$-algèbre formellement lisse
relativement à $K_\alpha$, pour les topologies discrètes (19.9.2, (iii)).
Mais comme $A$ est une $A_\alpha$-algèbre finie, on a
$A\otimes_{A_\alpha}K_\alpha=K$~; comme $\widehat A_\mathfrak p$ est un
anneau de fractions de $\widehat A\otimes_A K$, on en conclut (19.9.2,
(iv)) que $\widehat A_\mathfrak p$ est une $K$-algèbre formellement lisse
relativement à $K_\alpha$, pour la topologie discrète, et \emph{a fortiori}
pour sa topologie $\mathfrak p$-adique ((19.3.8) et (19.9.5)). Mais cela est
précisément la condition b) de (22.2.8), donc la proposition est démontrée
dans ce cas pour $p>0$. Lorsque $p=0$, le corps résiduel de
$\widehat A_\mathfrak p$ est séparable sur $K$, et comme
$\widehat A_\mathfrak p$ est un anneau régulier, c'est une $K$-algèbre
formellement lisse en vertu de (19.6.4).

II) Cas général. Soit $\mathfrak s$ l'idéal premier $\mathfrak n\cap A$ de
$A$, et posons $S=A-\mathfrak s$, de sorte que $S^{-1}A=A_\mathfrak s$, et
l'on a $C=(S^{-1}B)_{S^{-1}\mathfrak n}$, où cette fois
$S^{-1}\mathfrak n$ contient l'idéal maximal
$\mathfrak sA_\mathfrak s$ de $A_\mathfrak s$. Comme $A_\mathfrak s$ est
régulier (17.3.2) et $S^{-1}B$ une $A_\mathfrak s$-algèbre finie intègre
contenant $A_\mathfrak s$, tout revient à vérifier la condition (ii) pour
$A_\mathfrak s$ lorsque $p>0$~: or, posons
$\mathfrak s_\alpha=\mathfrak s\cap A_\alpha$ et considérons l'anneau local
$(A_\alpha)_{\mathfrak s_\alpha}$. Pour tout $t\notin\mathfrak s$, on a par
hypothèse, en posant $q=p^{h(\alpha)}$, $t^q\in A_\alpha$ et
$t^q\notin\mathfrak s$, donc $t^q\notin\mathfrak s_\alpha$~; si l'on pose
$S_\alpha=A_\alpha-\mathfrak s_\alpha$, on voit donc que l'on a
$A_\mathfrak s=S_\alpha^{-1}A$, et l'hypothèse entraîne que
$A_\mathfrak s$ est une $(A_\alpha)_{\mathfrak s_\alpha}$-algèbre finie~;
en outre, $K_\alpha$ est aussi le corps des fractions de
$(A_\alpha)_{\mathfrak s_\alpha}$, ce qui achève la démonstration.
\end{proof}

\begin{theorem}[22.3.3]
\label{0.22.3.3-fr}
Soient $A$ un anneau local noethérien complet, $\mathfrak p$ un idéal
premier de $A$, $B$ l'anneau localisé $A_\mathfrak p$, $\mathfrak q$ un
idéal premier du complété $\widehat B$, $\mathfrak r=\mathfrak q\cap B$.
Alors l'anneau localisé $\widehat B_\mathfrak q$ de $\widehat B$ est une
$B_\mathfrak r$-algèbre formellement lisse pour les topologies préadiques.
\end{theorem}

\begin{proof}
L'idéal premier $\mathfrak r$ de $B$ est de la forme
$\mathfrak n_\mathfrak p$, où $\mathfrak n\subset\mathfrak p$ est un idéal
premier de $A$, et l'on a $B_\mathfrak r=A_\mathfrak n$~; soit $k$ le corps
résiduel de $B_\mathfrak r$, qui est aussi le corps des fractions de l'anneau
local intègre et complet $A'=A/\mathfrak n$. Comme $\widehat B$ est un
$B$-module plat $(\mathbf 0_\mathrm I,~7.3.5)$,
$\widehat B_\mathfrak q$ est un $B_\mathfrak r$-module plat
$(\mathbf 0_\mathrm I,~6.3.2)$, et pour prouver que
$\widehat B_\mathfrak q$ est une $B_\mathfrak r$-algèbre formellement
lisse, il suffit de prouver que
\[
  \widehat B_\mathfrak q\otimes_{B_\mathfrak r}k
  =\widehat B_\mathfrak q/\mathfrak r\widehat B_\mathfrak q
\]
est une $k$-algèbre formellement lisse (19.7.1). Or
$\widehat B_\mathfrak q/\mathfrak r\widehat B_\mathfrak q
=(\widehat B/\mathfrak r\widehat B)_{\mathfrak q'}$, où $\mathfrak q'$ est
l'image canonique de $\mathfrak q$ dans
$\widehat B/\mathfrak r\widehat B$~; on a par définition
$\mathfrak q'\cap(B/\mathfrak r)=0$ et
$B/\mathfrak r=A_{\mathfrak p'}$, où
$\mathfrak p'=\mathfrak p/\mathfrak n$~; en outre
$\widehat B/\mathfrak r\widehat B$ est le complété de $B/\mathfrak r$.

On est donc ramené à démontrer le corollaire suivant~:
\end{proof}

\begin{corollary}[22.3.4]
\label{0.22.3.4-fr}
Soient $A$ un anneau local noethérien complet et intègre, $K$ son corps des
fractions, $\mathfrak p$ un idéal premier de $A$, $B$ l'anneau localisé
$A_\mathfrak p$. Alors, pour tout idéal premier $\mathfrak q$

